\section{Conclusion}

We present Hestia, voxel-face-aware hierarchical next-best-view acquisition for efficient 3D reconstruction. 
Hestia addresses the high-dimensional action space by separately predicting look-at points and camera positions.
Treating voxels as cubes enables more comprehensive capture, improving coverage ratios. 
The close-greedy design mitigates spurious correlations, ensuring efficient policy learning. 
Trained on a more diverse dataset, Hestia is robust across varied object-centric scenes.
Evaluations on three datasets validate that Hestia’s improvements are not marginal. 
Finally, the integration into a real-world drone system highlights its feasibility. 
As discussed in the limitations and future steps section (\cref{asec:limit}), one important step is extending to multi-agent settings to further improve efficiency.

\section{Acknowledgements}

This work was supported in part by the Australian Research Council (ARC) under discovery grant DP220100803 and DP250103612 and ARC Research Hub for Human-Robot Teaming for Sustainable and Resilient Construction (ITRH) grant IH240100016, and Australian National Health and Medical Research Council (NHMRC) Ideas Grant APP2021183. Research was also sponsored in part by the Australia Advanced Strategic Capabilities Accelerator (ASCA) under Contract No. P18-650825 and ASCA EDT DA ID12994, and the Australian Defence Science Technology Group (DSTG) under Agreement No: 12549.
We thank Yu-Lun Liu for submission guidance, Yi-Shan Hung for proofreading and demo narration, Meredith Porte for additional narration, and Xiao Chen for GenNBV~\cite{chen2024gennbv} discussions.
%We thank Yu-Lun Liu for insightful discussions on project direction and Yi-Shan Hung for assistance with paper proofreading and demo video narration.
%\clearpage